\[ %%%%%%%%%%%%%%%%%%%%%%%%%%%%%%%%%%%%%%%%%%%%%%%%%%%%%%%%%%%%%%%%%%%%%%%%%%%%%%%% \section{Experimental Methodology} \label{sec:experiments} %%%%%%%%%%%%%%%%%%%%%%%%%%%%%%%%%%%%%%%%%%%%%%%%%%%%%%%%%%%%%%%%%%%%%%%%%%%%%%%% This section describes the experimental methodology designed to evaluate the capabilities of PHYSGEN for long-horizon prediction of nonlinear dynamical systems. The experimental framework is designed to answer several fundamental questions: \begin{itemize} \item Can physics-guided graph networks improve long-horizon forecasting compared with purely data-driven approaches? \item Does explicit integration of physical constraints improve trajectory stability? \item Can PHYSGEN maintain physically consistent predictions under recursive rollout? \item How effectively can the framework generalize across different classes of dynamical systems? \end{itemize} %%%%%%%%%%%%%%%%%%%%%%%%%%%%%%%%%%%%%%%%%%%%%%%%%%%%%%%%%%%%%%%%%%%%%%%%%%%%%%%% \subsection{Research Questions} \label{sec:research_questions} %%%%%%%%%%%%%%%%%%%%%%%%%%%%%%%%%%%%%%%%%%%%%%%%%%%%%%%%%%%%%%%%%%%%%%%%%%%%%%%% The experimental evaluation is organized around five research questions (RQs). %%%%%%%%%%%%%%%%%%%%%%%%%%%%%%%%%%%%%%%%%%%%%%%%%%%%%%%%%%%%%%%%%%%%%%%%%%%%%%%% \subsubsection{RQ1: Long-Horizon Prediction Accuracy} %%%%%%%%%%%%%%%%%%%%%%%%%%%%%%%%%%%%%%%%%%%%%%%%%%%%%%%%%%%%%%%%%%%%%%%%%%%%%%%% \textbf{RQ1:} Can PHYSGEN achieve improved accuracy for long-horizon prediction compared with existing neural dynamical models? The objective is to evaluate whether physics-guided graph evolution provides advantages over conventional forecasting approaches when prediction horizons increase. The evaluation focuses on: \begin{equation} \mathcal{E}_{prediction} = \| X_{t+k} - \hat{X}_{t+k} \|. \end{equation} where $k$ represents the forecasting horizon. %%%%%%%%%%%%%%%%%%%%%%%%%%%%%%%%%%%%%%%%%%%%%%%%%%%%%%%%%%%%%%%%%%%%%%%%%%%%%%%% \subsubsection{RQ2: Physical Consistency} %%%%%%%%%%%%%%%%%%%%%%%%%%%%%%%%%%%%%%%%%%%%%%%%%%%%%%%%%%%%%%%%%%%%%%%%%%%%%%%% \textbf{RQ2:} Does PHYSGEN preserve physical constraints during autonomous prediction? This question evaluates whether the generated trajectories remain inside the physically admissible state space: \begin{equation} \hat{X}_t \in \mathcal{M}_{phys}. \end{equation} The analysis includes: \begin{itemize} \item conservation error; \item invariant violation; \item boundary condition consistency. \end{itemize} %%%%%%%%%%%%%%%%%%%%%%%%%%%%%%%%%%%%%%%%%%%%%%%%%%%%%%%%%%%%%%%%%%%%%%%%%%%%%%%% \subsubsection{RQ3: Stability of Long-Term Rollout} %%%%%%%%%%%%%%%%%%%%%%%%%%%%%%%%%%%%%%%%%%%%%%%%%%%%%%%%%%%%%%%%%%%%%%%%%%%%%%%% \textbf{RQ3:} Can PHYSGEN reduce error accumulation during recursive multi-step forecasting? The stability evaluation measures trajectory divergence: \begin{equation} D_k = \| X_{t+k} - \hat{X}_{t+k} \|. \end{equation} The objective is to analyze error growth as a function of prediction horizon. %%%%%%%%%%%%%%%%%%%%%%%%%%%%%%%%%%%%%%%%%%%%%%%%%%%%%%%%%%%%%%%%%%%%%%%%%%%%%%%% \subsubsection{RQ4: Contribution of Individual Components} %%%%%%%%%%%%%%%%%%%%%%%%%%%%%%%%%%%%%%%%%%%%%%%%%%%%%%%%%%%%%%%%%%%%%%%%%%%%%%%% \textbf{RQ4:} What is the contribution of individual PHYSGEN components? Ablation studies are performed by removing: \begin{itemize} \item physics-guided message passing; \item stability control module; \item multi-level physical constraints; \item graph adaptation mechanism. \end{itemize} The goal is to quantify the contribution of each architectural component. %%%%%%%%%%%%%%%%%%%%%%%%%%%%%%%%%%%%%%%%%%%%%%%%%%%%%%%%%%%%%%%%%%%%%%%%%%%%%%%% \subsubsection{RQ5: Cross-System Generalization} %%%%%%%%%%%%%%%%%%%%%%%%%%%%%%%%%%%%%%%%%%%%%%%%%%%%%%%%%%%%%%%%%%%%%%%%%%%%%%%% \textbf{RQ5:} Can PHYSGEN generalize across different nonlinear physical systems? The evaluation considers transfer between different dynamical regimes: \begin{equation} \mathcal{D}_i \rightarrow \mathcal{D}_j, \quad i\neq j. \end{equation} This experiment investigates whether the learned physical representations capture general dynamical principles rather than dataset-specific correlations. %%%%%%%%%%%%%%%%%%%%%%%%%%%%%%%%%%%%%%%%%%%%%%%%%%%%%%%%%%%%%%%%%%%%%%%%%%%%%%%% \subsubsection{Experimental Hypotheses} %%%%%%%%%%%%%%%%%%%%%%%%%%%%%%%%%%%%%%%%%%%%%%%%%%%%%%%%%%%%%%%%%%%%%%%%%%%%%%%% Based on the PHYSGEN design principles, we formulate the following hypotheses: \paragraph{H1.} Physics-guided graph evolution improves long-horizon prediction accuracy compared with purely data-driven baselines. \paragraph{H2.} Multi-level physical constraints reduce violations of conservation laws and physical invariants. \paragraph{H3.} Adaptive stability control decreases error growth during recursive rollout. \paragraph{H4.} The learned physical representations improve transfer across different dynamical systems. %%%%%%%%%%%%%%%%%%%%%%%%%%%%%%%%%%%%%%%%%%%%%%%%%%%%%%%%%%%%%%%%%%%%%%%%%%%%%%%% \subsubsection{Summary} %%%%%%%%%%%%%%%%%%%%%%%%%%%%%%%%%%%%%%%%%%%%%%%%%%%%%%%%%%%%%%%%%%%%%%%%%%%%%%%% The experimental methodology is designed to evaluate PHYSGEN from multiple perspectives, including prediction accuracy, physical consistency, stability, component contribution, and cross-domain generalization. Together, these experiments provide a comprehensive evaluation framework for physics-guided graph learning of nonlinear dynamical systems. \] 
